\section{Nearly full-rate cancellations}
\label{can:section}

The first estimate concerns two long products which expand rapidly when
viewed from their common endpoint, although their concatenation has little
growth.  Such a configuration forces a linearized solution to be small at
both distant endpoints.  We use this solution to locate many pairs of
comparable expansion scales on the two sides of the common endpoint.

\begin{theorem}[Cancellation estimate]\label{can:theorem}
There are constants $\delta>0$, $C<\infty$, $\gamma>0$, and $M_0<\infty$
such that, for every $M\ge M_0$, every $i\in\Z$, and every $n\in\N$,
\begin{equation}\label{can:event-bound}
\PP\left\{
\begin{array}{l}
\min\bigl(\log_M\norm{T_{i,i-n}},
          \log_M\norm{T_{i,i+n}}\bigr)\ge(1-\delta)n,\\[2pt]
\log_M\norm{T_{i-n,i+n}}\le\delta n
\end{array}
\right\}
\le C\exp(-\gamma n).
\end{equation}
All four constants are independent of the center, the length, and the
parameter $M$ in the indicated range.
\end{theorem}

Throughout this section and the next, fix
\begin{equation}\label{can:c}
c=\frac{99}{100}.
\end{equation}
Constants used before the final choice of $\delta$ are independent of
$\delta$, $n$, and $M$, unless a dependence is stated explicitly.

\subsection{A solution small at both endpoints}

By stationarity it is enough to consider $i=0$.  Suppose the event in
\cref{can:event-bound} occurs.  Write
\[
D_- = \norm{T_{0,-n}},\qquad D_+ = \norm{T_{0,n}},
\]
and let $a_-$ and $a_+$ be unit most-expanded input directions for these
two products.  If $b$ is a unit vector perpendicular to $a_-$, then
\[
\norm{T_{0,-n}b}=D_-^{-1},\qquad
\norm{T_{0,n}b}\ge D_+\abs{a_-\wedge a_+}.
\]
The first equality uses determinant one: the singular values of either
product are its norm and the reciprocal of its norm.  Applying the full
product to the vector $T_{0,-n}b$ gives
\[
\norm{T_{-n,n}}
\ge D_-D_+\abs{a_-\wedge a_+},
\qquad
\abs{a_-\wedge a_+}\le M^{-(2-3\delta)n}.
\]
Resolving $b$ in the singular input directions of $T_{0,n}$ also gives the
upper bound
\[
\norm{T_{0,n}b}
\le D_+^{-1}+D_+\abs{a_-\wedge a_+}
\le 2M^{-(1-3\delta)n}.
\]
The left endpoint has the smaller bound $D_-^{-1}$.  Thus the solution
of the linearized recurrence with central pair $b$ is small at both
endpoints.  One of its entries at times $0$ and $-1$ has absolute value
at least $1/\sqrt2$.  Choose that time $e\in\{0,-1\}$ and divide the
whole solution by its value there.  The normalized solution satisfies
\begin{equation}\label{can:small-solution}
u_e=1,\qquad \norm{(u_0,u_{-1})}\le\sqrt2,\qquad
\norm{(u_n,u_{n-1})},\ \norm{(u_{-n},u_{-n-1})}
\le 3M^{-(1-3\delta)n}.
\end{equation}
This construction will be used only as an existence statement; no
measurable selection of a singular direction is required.

\subsection{Outward paths and good indices}

Fix $e$ and $q_e$.  Use $\theta=q_{e-1}$ as the other circle coordinate,
so that $q_{e+1}=\phi(q_e)-\theta$ modulo one.  The two outward paths are
\[
q_j^-=q_{e-j},\qquad q_j^+=q_{e+j}\qquad(j\ge0).
\]
On either path the positions and every linearized solution obey the same
second-order recurrence.  In calculations involving just one path, omit
its sign and write $q_j,u_j,v_j=v(q_j)$, with $u_0=1$.  Define
\begin{equation}\label{can:t-definition}
t_0=0,\qquad t_1=1,\qquad t_{j+1}=v_jt_j-t_{j-1}.
\end{equation}
For local real lifts of the positions,
\[
\partial_\theta q_j^-=t_j^-,\qquad
\partial_\theta q_j^+=-t_j^+.
\]
The sign difference comes solely from the initial positions at index one.
It does not change the recurrence or its Wronskian identity:
\begin{equation}\label{can:wronskian}
u_{j-1}t_j-u_jt_{j-1}=1\qquad(j\ge1).
\end{equation}
Indeed the expression on the left is one at $j=1$ and the recurrence
preserves it.  A common normalized solution on the two paths has
\begin{equation}\label{can:center-equation}
u_1^-+u_1^+=v(q_e).
\end{equation}
Conversely, two outward solutions with value one at index zero and
satisfying \cref{can:center-equation} assemble into a solution of the
original recurrence across the center.

On each path the two small endpoint entries in
\cref{can:small-solution} are indexed by $N-1,N$. More explicitly, the
forward and backward lengths are $N_+=n-e$ and $N_-=n+e+1$,
respectively: one is $n$ and the other is $n+1$. Define pair-growth
logarithms
\[
\tau_j=\log_M\max(\abs{t_j},\abs{t_{j-1}}),\qquad
\upsilon_j=\log_M\max(\abs{u_j},\abs{u_{j-1}})
\quad(1\le j\le N).
\]
Both pairs are nonzero, so these logarithms are defined.  The bounds for
one step and its inverse in the maximum norm give
\begin{equation}\label{can:log-increments}
\abs{\tau_{j+1}-\tau_j}\le1,\qquad
\abs{\upsilon_{j+1}-\upsilon_j}\le1.
\end{equation}

\begin{definition}\label{can:good-definition}
An index $2\le j\le N-2$ is \emph{good} for the selected normalized
solution if
\begin{equation}\label{can:good}
\begin{aligned}
\tau_j-\tau_l&\ge c(j-l)&& (1\le l<j),\\
\tau_{j+1}-\tau_j&\ge c,\\
\upsilon_j-\upsilon_{j+1}&\ge c,
&\upsilon_{j+1}-\upsilon_{j+2}&\ge c.
\end{aligned}
\end{equation}
For a path extending beyond a selected index, the same inequalities
define goodness without reference to a particular terminal length $N$.
\end{definition}

\begin{lemma}\label{can:many-good}
There is a constant $A$ such that, on each outward path associated to
\cref{can:small-solution}, the number of indices in $\{2,\ldots,N-2\}$
which are not good is at most
\begin{equation}\label{can:bad-count}
\frac{A(\delta n+1)}{1-c}.
\end{equation}
The estimate holds uniformly for $M\ge e$.
\end{lemma}

\begin{proof}
At index $N$, \cref{can:wronskian} and the endpoint bound imply
\[
1\le 2\max(\abs{u_N},\abs{u_{N-1}})
          \max(\abs{t_N},\abs{t_{N-1}}),
\]
and hence
\[
\tau_N\ge(1-3\delta)n-\log_M6,
\qquad
\upsilon_N\le-(1-3\delta)n+\log_M3.
\]
Since $\tau_1=0$, $\upsilon_1\ge0$, and $N-1\le n$, the total
shortfalls of the paths $\tau$ and $-\upsilon$ satisfy
\begin{align*}
\sum_{j=1}^{N-1}\bigl[1-(\tau_{j+1}-\tau_j)\bigr]
&\le3\delta n+\log6,\\
\sum_{j=1}^{N-1}\bigl[1-(\upsilon_j-\upsilon_{j+1})\bigr]
&\le3\delta n+\log3.
\end{align*}
Every summand is nonnegative by \cref{can:log-increments}.  Apply
\cref{setup:shortfall-count} to $\tau$ to count indices where the first
line of \cref{can:good} fails.  An individual increment smaller than $c$
costs at least $1-c$ from the relevant shortfall sum.  This counts
failures of the second line; applied to $-\upsilon$, it counts failures
of either of the two inequalities in the last line, with a factor of at
most two because a single increment can be tested at two neighboring
indices.  Summing these bounds proves \cref{can:bad-count}.
\end{proof}

At a good index let $L_j=\abs{t_j}$.  The first line of
\cref{can:good} with $l=j-1$ makes $t_j$ the dominant entry of its pair,
and the next increase makes $t_{j+1}$ dominant at the next pair.
Consequently
\begin{equation}\label{can:good-t}
\begin{gathered}
L_j=M^{\tau_j},\qquad
\abs{t_l}\le L_jM^{-c(j-l)}\quad(1\le l<j),\\
\abs{t_{j+1}}\ge M^cL_j.
\end{gathered}
\end{equation}
The two decreases for the $u$-pairs similarly give
\[
\abs{u_{j+1}}\le M^{-c}\abs{u_j}
\le M^{-2c}\abs{u_{j-1}}.
\]
Using \cref{can:wronskian} and
$\abs{t_{j-1}/t_j}\le M^{-c}$, we obtain, after increasing a fixed
lower bound for $M$,
\begin{equation}\label{can:good-u}
\abs{u_{j-1}}L_j
\le 1+M^{-2c}\abs{u_{j-1}}L_j,
\qquad \abs{u_{j-1}}\le\frac2{L_j}.
\end{equation}

For $t_j\ne0$, let $R_j$ be the initial value at index one of the
solution having value one at index zero and value zero at index $j$.
This solution exists and is unique: changing its initial value at index
one by $a$ changes its value at index $j$ by $at_j$.  Its value at
index $j-1$ is $1/t_j$, by \cref{can:wronskian}. Its value at index
$j+1$ is therefore $-1/t_j$, which proves
\[
R_{j+1}-R_j=\frac1{t_jt_{j+1}}.
\]
For any normalized solution, varying its initial value at index one also
gives $u_1-R_{j+1}=u_{j+1}/t_{j+1}$.  At a good index,
\cref{can:good-t,can:good-u} thus yield
\begin{equation}\label{can:boundary-expansion}
\begin{aligned}
u_1
&=R_{j+1}+\frac{u_{j+1}}{t_{j+1}}\\
&=R_j+\frac1{t_j^2(v_j-t_{j-1}/t_j)}
  +O\left(\frac{M^{-3c}}{L_j^2}\right).
\end{aligned}
\end{equation}
The implied constant is absolute and does not depend on the index.

\subsection{Matching the two expansion scales}

\Cref{can:wing-figure} depicts the two outward paths and one selected
pair with comparable logarithmic derivative scales, as in the following lemma.

\begin{lemma}\label{can:matching}
There are fixed constants $\delta_0>0$, $c_1>0$, $C_1<\infty$, and
$n_*<\infty$ with the following property.  If $0<\delta\le\delta_0$
is fixed, increasing $n_*$ if necessary, every configuration in
\cref{can:small-solution} with $n\ge n_*$ has at least $c_1n$ pairs
$(j_m^+,j_m^-)$ of good indices such that both index sequences increase
strictly and
\begin{equation}\label{can:matching-condition}
\abs{\tau_{j_m^+}^+-\tau_{j_m^-}^-}\le\frac35.
\end{equation}
The collection of all possible increasing index-pair patterns for a
given $n$ has cardinality at most $\exp(C_1n)$.
\end{lemma}

\begin{proof}
Choose a small fixed $\beta>0$, for example $\beta=10^{-3}$.  By
\cref{can:bad-count}, choosing $\delta_0$ sufficiently small and then
$n_*$ sufficiently large makes the bad-index count on each path at most
$\beta n$.  The total-shortfall estimate also gives, at every index,
\[
j-1-3\delta n-\log6\le\tau_j\le j-1.
\]
Shrink $\delta_0$ once more and increase $n_*$ so that, for
$n/3\le j\le n/2$, the corresponding $\tau_j$ lies in
$[n/4,3n/4]$, while on the minus path
\[
\tau_{N-2}^-\ge\tau_N^--2>3n/4+1.
\]
There are at least $n/6-2-\beta n$ good plus indices in this range.

For each such index put $a=\tau_j^+$.  Starting from $\tau_1^-=0$,
let $k$ be the first minus index with $\tau_k^-\ge a$.  We have
$\tau_{k-1}^-<a\le\tau_k^-$ and
$\tau_k^--\tau_{k-1}^-\le1$.  Thus at least one of these two values
is within $1/2$ of $a$.  The chosen index is in
$\{2,\ldots,N-2\}$: the large lower bound on $a$ excludes the first
two indices, and the displayed lower bound at $N-2$ provides the upper
bound.  If the chosen minus index is bad, discard this plus index.

Good plus indices $j<j'$ have
$\tau_{j'}^+-\tau_j^+\ge c(j'-j)$.  At most two such logarithms can
belong to an interval of length one, since $2c>1$.  A fixed bad minus
index can therefore be responsible for at most two discarded plus
indices.  At least $n/6-2-3\beta n$ candidates remain.

Thin these candidates in increasing plus order by retaining the first
one, then retaining the next whose plus logarithm is at least two larger
than the last retained logarithm, and continuing.  Since successive
candidate logarithms are separated by at least $c$, each retention
removes at most three candidates.  A fixed positive fraction of $n$
remains; for example, decreasing the fraction if necessary, one may fix
$c_1=10^{-2}$.  The matched minus logarithms increase strictly, because
each is within $1/2$ of the corresponding plus logarithm and the retained
plus logarithms have gaps at least two.  For two good minus indices,
increasing index implies strictly increasing logarithm by
\cref{can:good}.  Strict increase of their logarithms therefore forces
strict increase of their indices as well.  The matching error is at most
$1/2$, which implies \cref{can:matching-condition}.

Each path contains at most $n+1$ available indices.  An increasing paired
pattern is determined by a subset of the indices on each path, with the
two subsets of equal cardinality paired in order.  There are at most
$2^{2n+2}$ such choices.  Enlarging an absolute $C_1$ gives the stated
exponential bound.
\end{proof}

\begin{figure}[tb]
\centering
\begin{tikzpicture}[x=1cm,y=1cm,>=Stealth]
  \draw[->] (0,0)--(-5,0);
  \draw[->] (0,0)--(5,0);
  \foreach \x in {-4.6,-2,0,2.6,4.6}
    \fill (\x,0) circle (1.5pt);
  \node[above] at (0,0) {$q_e$};
  \node[above] at (-2,0) {$q_{e-j_m^-}$};
  \node[above] at (2.6,0) {$q_{e+j_m^+}$};
  \node[above] at (-4.6,0) {left endpoint};
  \node[above] at (4.6,0) {right endpoint};
  \draw[densely dashed] (-2,-.15)--(-2,-.65)--(2.6,-.65)--(2.6,-.15);
  \node[below,align=center] at (0,-.65)
    {a matched pair: $\abs{\log_M L_{j_m^-}^- -\log_M L_{j_m^+}^+}\le3/5$};
\end{tikzpicture}
\caption{The two paths are followed outward from a fixed central position.
Matched indices need not be equally far from the center; their derivatives
with respect to the other central coordinate have comparable logarithmic
size.}
\label{can:wing-figure}
\end{figure}

\subsection{The local input}

The analytic fact used at each matched pair is the following uniform
scarcity statement.  Its proof occupies \cref{loc:section}.

\begin{lemma}[Local scarcity]\label{can:local-scarcity}
For every $\eta>0$ there are $\eta'>0$ and $M_\eta<\infty$ such that
the following holds for $M\ge M_\eta$.  Fix $q_e$, a base coordinate
$\theta_*$, and indices $j_+,j_-\ge2$ on the two outward paths.
Suppose both indices are good at $\theta_*$ for a common normalized
solution, and their logarithms satisfy
\cref{can:matching-condition}.  Let
\[
L_+=\abs{t_{j_+}^+(\theta_*)},\qquad
L_-=\abs{t_{j_-}^-(\theta_*)},\qquad L=\max(L_+,L_-).
\]
Let $\mathcal S\subset[-1,1]$ consist of the $h$ for which, at
$\theta=\theta_*+h/L$, the same two indices are good for some common
solution with $u_e=1$.  The solution is allowed to depend on $h$, and
neither endpoint smallness nor a new matching condition is required at
these other points.  Then
\begin{equation}\label{can:scarcity-bound}
\abs{\{x\in\R:\dist(x,\mathcal S)<\eta'\}}\le\eta.
\end{equation}
Here and below the size of a subset of the line or circle denotes its
Lebesgue measure.  The constants are uniform over the central position,
the base coordinate, and the selected indices.
\end{lemma}

\subsection{From local scarcity to exponential decay}

\begin{proof}[Proof of \cref{can:theorem}]
Fix once and for all a sufficiently small $\delta>0$ allowed by
\cref{can:matching}.  Its constants $c_1,C_1,n_*$ are now fixed.
Choose $\eta>0$ so small that
\begin{equation}\label{can:eta-choice}
5\eta<1,\qquad
c_1\log\frac1{5\eta}>C_1+1.
\end{equation}
Obtain $\eta'$ and $M_\eta$ from \cref{can:local-scarcity}.  Increase
a fixed lower bound $M_0$ so that all estimates above hold, that
$M_0\ge M_\eta$, and that
\begin{equation}\label{can:radius-choice}
M_0^{-c}<\frac14,\qquad M_0^{-c}/5<\eta'.
\end{equation}
These choices are independent of $n$, of the center, and of $M\ge M_0$.

Fix $e\in\{0,-1\}$ and $q_e$.  It is enough to consider patterns
containing exactly $r=\lceil c_1n\rceil$ pairs, since a longer pattern
can be truncated.  For a fixed such increasing pattern, let
$\mathcal F$ be the set of $\theta$ admitting one common normalized
solution with all of the following properties: the original central
pair has norm at most $\sqrt2$; both endpoint pairs satisfy
\cref{can:small-solution}; and every pair of selected indices is good
and satisfies \cref{can:matching-condition}.  We emphasize that one
solution must satisfy all conditions of a pattern at once.

The set $\mathcal F$ is measurable, in fact compact.  To see this without
selecting solutions, parameterize the original central pair by its entry
other than the fixed entry $u_e=1$.  The norm bound puts that parameter in
$[-1,1]$.  All later solution entries and all recurrence coefficients
depend continuously on this parameter and on the circle coordinate
$\theta$.  Each pair entering a logarithm is nonzero, by invertibility
of the recurrence.  The endpoint, goodness, and matching conditions are
closed inequalities on the compact product of the circle with
$[-1,1]$.  Their projection is therefore compact.

For $\theta\in\mathcal F$ and $1\le m\le r$, define
\[
L_m(\theta)=\max\left(
\abs{t_{j_m^+}^+(\theta)},\abs{t_{j_m^-}^-(\theta)}\right),
\qquad s_m(\theta)=L_m(\theta)^{-1},
\]
and let
\[
V_m=\bigcup_{\theta\in\mathcal F}
 B\bigl(\theta,s_m(\theta)/10\bigr)
\]
be an open subset of the circle.  Goodness at the later of two selected
indices gives, separately on both paths,
\[
\abs{t_{j_{m+1}^\pm}^\pm(\theta)}
\ge M^{c(j_{m+1}^\pm-j_m^\pm)}
      \abs{t_{j_m^\pm}^\pm(\theta)}.
\]
Since both index sequences increase strictly,
\begin{equation}\label{can:radius-decay}
s_{m+1}(\theta)\le M^{-c}s_m(\theta),
\qquad s_m(\theta)\le M^{-c}.
\end{equation}
For fixed $M,n$, these positive radii also have a positive lower bound:
the one-step norm bound gives $\abs{t_j}\le M^{j-1}\le M^n$ at all
indices in question, so $s_m\ge M^{-n}$.

We claim
\begin{equation}\label{can:cover-contraction}
\abs{V_{m+1}}\le5\eta\abs{V_m}\qquad(1\le m<r).
\end{equation}
Select disjoint balls from the family defining $V_m$ as follows.  At each
stage select a remaining ball whose radius is at least half the
supremum of the remaining radii, and remove all remaining balls meeting
it.  Assign each removed ball to this selected ball.  All selected balls
are disjoint and have a fixed positive lower radius, so only finitely
many can be selected on the circle; the procedure must exhaust the
family.  Every assigned ball meets its selected ball and has radius at
most twice that of the selected ball.

Consider one selected center $\theta_*$ and write
$s=s_m(\theta_*)$.  The centers of its assigned balls are at circle
distance at most $3s/10$ from $\theta_*$.  The window of radius $s$ is
an injective lift interval by \cref{can:radius-choice,can:radius-decay}.
All assigned centers belong to $\mathcal F$, so they are successful for
the same selected pair of indices in the local test based at
$\theta_*$.  A next-level ball centered at an assigned $\theta$ has
radius at most
\[
\frac{s_{m+1}(\theta)}{10}
\le \frac{M^{-c}s_m(\theta)}{10}
\le \frac{M^{-c}s}{5}.
\]
After rescaling the lifted coordinate by $1/s$, the union of all these
next-level balls is contained in the $\eta'$-neighborhood of the
successful set in \cref{can:local-scarcity}, by
\cref{can:radius-choice}.  It therefore has unscaled length at most
$\eta s$.  Summing over selected balls proves
\[
\abs{V_{m+1}}\le\eta\sum_{\text{selected}}s
=5\eta\sum_{\text{selected}}
  \abs{B(\theta_*,s/10)}
\le5\eta\abs{V_m},
\]
as claimed.  The ball selected at a center is included in its own
assigned group, so no next-level balls have been omitted.

Since $\mathcal F\subset V_m$ for every $m$ and $\abs{V_1}\le1$,
iteration gives
\[
\abs{\mathcal F}\le(5\eta)^{r-1}
\le(5\eta)^{-1}
  \exp\left(-c_1n\log\frac1{5\eta}\right).
\]
There are at most $\exp(C_1n)$ patterns.  By
\cref{can:eta-choice}, their union has circle measure at most
$C\exp(-\gamma n)$ for a fixed $\gamma>0$, uniformly in the fixed
coordinate $q_e$.  The coordinate pair $(q_e,q_{e-1})$ has normalized
product Lebesgue law.  Integrating the conditional bound over $q_e$ and
summing the two choices of $e$ therefore gives the same form of bound
for the original event, for every $n\ge n_*$.  Every point of that event
is covered, by \cref{can:small-solution,can:matching}.  Finally enlarge
$C$ to cover the finitely many positive integers $n<n_*$, whose event
probabilities are at most one.  Stationarity restores every center $i$.
\end{proof}
